\chapter{แนวคิด ทฤษฎี และงานวิจัยที่เกี่ยวข้อง}

\section{บริบทภาคการเกษตรของประเทศไทยและนโยบายเกษตรกรอัจฉริยะ}

ภาคการเกษตรเป็นหนึ่งในภาคเศรษฐกิจหลักของประเทศไทย โดยมีแรงงานมากกว่าหนึ่งในสามของกำลังแรงงานทั้งหมด แต่สร้างมูลค่าเพียงประมาณร้อยละ 10 ของผลิตภัณฑ์มวลรวมในประเทศ ซึ่งสะท้อนถึงปัญหาผลิตภาพที่ต่ำ \parencite{un2020}.

เกษตรกรไทยต้องเผชิญกับปัญหาหลายประการ เช่น ความผันผวนของสภาพภูมิอากาศ การขาดแคลนทรัพยากร และข้อจำกัดในการเข้าถึงเทคโนโลยี \parencite{chatarat2019}.

เพื่อตอบสนองต่อปัญหาดังกล่าว ภาครัฐได้ดำเนินนโยบาย “เกษตรกรอัจฉริยะ (Smart Farmer)” เพื่อพัฒนาศักยภาพของเกษตรกรให้สามารถใช้ข้อมูลและเทคโนโลยีในการตัดสินใจ \parencite{nation2023}.

---

\section{แนวคิดและองค์ประกอบของ Smart Farmer}

แนวคิด Smart Farmer มุ่งเน้นให้เกษตรกรมีความสามารถในการบริหารจัดการการผลิต โดยใช้ข้อมูลและเทคโนโลยีสนับสนุนการตัดสินใจ

องค์ประกอบสำคัญ ได้แก่:

\begin{itemize}
\item \textbf{ด้านองค์ความรู้} ความเข้าใจในกระบวนการผลิต
\item \textbf{ด้านการจัดการ} การวางแผนการผลิตและการตลาด
\item \textbf{ด้านข้อมูล} การใช้ข้อมูลในการตัดสินใจ
\item \textbf{ด้านเทคโนโลยี} การนำเทคโนโลยีมาใช้เพิ่มผลิตภาพ
\end{itemize}

---

\section{แนวคิดด้านเทคโนโลยีการเรียนรู้ของเครื่อง (Machine Learning)}

Machine Learning เป็นเทคนิคทางปัญญาประดิษฐ์ที่ช่วยให้ระบบสามารถเรียนรู้จากข้อมูลและปรับปรุงผลลัพธ์ได้เอง \parencite{aggarwal2018}.

Machine Learning แบ่งออกเป็น 3 ประเภท ได้แก่:

\begin{itemize}
\item Supervised Learning
\item Unsupervised Learning
\item Reinforcement Learning
\end{itemize}

ในงานวิจัยนี้เลือกใช้ Supervised Learning เนื่องจากสามารถใช้ในการจำแนกและพยากรณ์ข้อมูลได้อย่างแม่นยำ

---

\section{การประยุกต์ใช้ Machine Learning ในภาคการเกษตร}

Machine Learning ถูกนำมาใช้ในภาคการเกษตร เช่น การพยากรณ์ผลผลิต การวิเคราะห์พฤติกรรมเกษตรกร และการจำแนกกลุ่มเป้าหมาย \parencite{jinjarak2007}.

นอกจากนี้ Machine Learning ยังสามารถช่วยออกแบบระบบฝึกอบรมให้เหมาะสมกับเกษตรกรแต่ละกลุ่มได้ \parencite{charoensukmongkol2015}.

---

\section{งานวิจัยที่เกี่ยวข้อง}

งานวิจัยก่อนหน้าพบว่า ปัจจัยด้านการศึกษา ทรัพยากร และการเข้าถึงข้อมูล มีผลต่อการเข้าร่วมโครงการเกษตร \parencite{abdallah2016}.

การศึกษาโดย \textcite{anang2020} พบว่าการเข้าร่วมโครงการส่งเสริมสามารถเพิ่มผลิตภาพได้ แต่ขึ้นอยู่กับบริบทของพื้นที่

---

\section{กรอบแนวคิดในการวิจัย (Conceptual Framework)}

จากการทบทวนวรรณกรรม สามารถกำหนดตัวแปรได้ดังนี้:

\begin{itemize}
\item ตัวแปรตาม: ทักษะของเกษตรกร
\item ตัวแปรอิสระ: การเข้าร่วมโครงการ และลักษณะส่วนบุคคล
\end{itemize}

ความสัมพันธ์ของตัวแปรถูกวิเคราะห์ด้วย Difference-in-Differences (DID) \parencite{wooldridge2009}.